% GENERATED FILE. Edit canonical lessons, then run npm run generate:books.
% canonical-source-hash: fnv1a64:db6cd2f3d2ad77fa
% canonical-lessons: BN-W04-pa, BN-W03-da, BN-W03-dha, BN-W04-cha, BN-W04-chha, BN-W04-achhi-read, BN-W04-ja

\chapter{The Roof, the Teeth, and the Grid Closed}
\label{ch:bn-roof-teeth-grid}
\hlchaptermodality{9}

\begin{chapteropening}
\textbf{By the end of this chapter:} \emph{I can write the five consonants that finish Bengali's stop grid, hear the breath that separates \bn{চ} from \bn{ছ}, read \bn{আছি} off the page, and say every consonant I own by where in the mouth it is made rather than by what it looks like.}

The chapter closes by reciting ten consonants back from their places of articulation rather than from their shapes, which is the difference between holding a grid and holding a list.
\end{chapteropening}

\section[pô]{\bn{প} — the consonant pa}

\label{lesson:BN-W04-pa}



\begin{quote}
Write \textbf{\bn{খ}} from memory, and say what separates it from \textbf{\bn{ক}}.
\end{quote}

\subsection*{Script}
\begin{quote}
\bn{প}
\end{quote}

\textbf{pô} — the \emph{p} of English \emph{spin}, lips shut, no puff of air after the release.

The place is the lips, which is where \textbf{\bn{ব}} was made too. So the two of them fill two cells of one square:

\begin{itemize}
\raggedright
  \item \textbf{\bn{প}} — lips, voice off
  \item \textbf{\bn{ব}} — lips, voice on
\end{itemize}

The English \emph{p} of \emph{pin} is the breathy one and belongs to a different letter, the same trap \textbf{\bn{ক}} and \textbf{\bn{খ}} laid a lesson ago. Palm in front of the mouth: \emph{spin} gives no air, and that is this letter.

With the lips named, the map is complete. Every consonant this book has taught is made at one of four places — \textbf{lips, teeth, roof, throat} — and every one of them is one of four cells at that place. A reader who has that grid is not memorising forty unrelated shapes. They are filling in a table.

With the mātrā:

\begin{quote}
\textbf{\bn{প}} + \textbf{\bn{া}} $\to$ \textbf{\bn{পা}}
\end{quote}

\subsection*{Your turn}
\begin{itemize}
\raggedright
  \item \emph{Write it:} \bn{প}
  \item \emph{Write it:} \bn{প} \bn{ব}
  \item \emph{Read it:} \bn{ভালো}
\end{itemize}

\begin{tcolorbox}[breakable,colback=teal!4,colframe=teal!35!black,title={Before you move on}]
Where is \textbf{\bn{প}} made? (\textbf{At the lips}.) Which other letter is made there? (\textbf{\bn{ব}}.) And what separates them? (\textbf{The voice}.)

Source: \href{https://www.unicode.org/charts/PDF/U0980.pdf}{Unicode Bengali chart}.
\end{tcolorbox}
\section[dô]{\bn{দ} — the consonant da}

\label{lesson:BN-W03-da}



\begin{quote}
Write \textbf{\bn{খ}} and \textbf{\bn{প}}, then write \textbf{\bn{ত}} and say where the tongue goes for it.
\end{quote}

\subsection*{Script}
\begin{quote}
\bn{দ}
\end{quote}

\textbf{dô} — the same tongue position as \textbf{\bn{ত}}, on the \textbf{back of the top front teeth}, with the voice switched on.

So the warning that came with \textbf{\bn{ত}} applies here without a word changed. English makes its \emph{d} on the ridge behind the teeth; Bengali has a separate letter for that position, and it is a different letter spelling different words. Slide the tongue forward onto enamel.

The pair to hold is not \emph{t} and \emph{d} as English thinks of them. It is:

\begin{itemize}
\raggedright
  \item \textbf{\bn{ত}} dental, voice off
  \item \textbf{\bn{দ}} dental, voice on
\end{itemize}

Two settings on one place of articulation. Bengali's consonant grid is built out of exactly these two switches — where the tongue is, and whether the voice is on — plus the breath from the last lesson. Three switches, and most of the inventory falls out of them.

With the mātrā:

\begin{quote}
\textbf{\bn{দ}} + \textbf{\bn{া}} $\to$ \textbf{\bn{দা}}
\end{quote}

\subsection*{Your turn}
\begin{itemize}
\raggedright
  \item \emph{Write it:} \bn{দ}
  \item \emph{Write it:} \bn{ত} \bn{দ}
  \item \emph{Read it:} \bn{ভালো} \bn{আবার}
\end{itemize}

\begin{tcolorbox}[breakable,colback=teal!4,colframe=teal!35!black,title={Before you move on}]
What is the same about \textbf{\bn{ত}} and \textbf{\bn{দ}}? (\textbf{Where the tongue goes} — the teeth.) What is different? (\textbf{The voice}.) And is an English \emph{d} the same letter? (\textbf{No} — it is made further back.)

Source: \href{https://www.unicode.org/charts/PDF/U0980.pdf}{Unicode Bengali chart}.
\end{tcolorbox}
\section[dhô]{\bn{ধ} — the consonant dha}

\label{lesson:BN-W03-dha}



\begin{quote}
Write \textbf{\bn{দ}} from memory, and say where the tongue touches for it.
\end{quote}

\subsection*{Script}
\begin{quote}
\bn{ধ}
\end{quote}

\textbf{dhô} — \textbf{\bn{দ}} with the breath let out through it. Voice on, tongue on the teeth, and air after the release.

Set the four dental letters side by side and the grid the last two lessons have been building is visible in one line:

\begin{itemize}
\raggedright
  \item \textbf{\bn{ত}} — voice off, no breath
  \item \textbf{\bn{দ}} — voice on, no breath
  \item \textbf{\bn{ধ}} — voice on, with breath
\end{itemize}

The fourth cell, voice off and with breath, is a letter this book has not needed yet. It exists, it sits between the first two in every dictionary, and it will arrive when a word calls for it.

\textbf{That square is the unit.} Every place of articulation in the script — teeth, ridge, palate, throat — gets the same four cells in the same order, and the order is over two thousand years old: the Sanskrit grammarians sorted the sounds by where they were made and wrote the list down, and every script descended from that tradition still lists them that way.

With the u-sign:

\begin{quote}
\textbf{\bn{ধ}} + \textbf{\bn{ু}} $\to$ \textbf{\bn{ধু}}
\end{quote}

\subsection*{Your turn}
Write these out:

\begin{itemize}
\raggedright
  \item \bn{ধ}
  \item \bn{ত} \bn{দ} \bn{ধ}
  \item \bn{দা} \bn{ধা}
\end{itemize}

\begin{tcolorbox}[breakable,colback=teal!4,colframe=teal!35!black,title={Before you move on}]
What two things separate \textbf{\bn{ধ}} from \textbf{\bn{ত}}? (\textbf{Voice} and \textbf{breath}.) How many cells does one place of articulation get? (\textbf{Four}.) And who fixed that order? (The \textbf{Sanskrit grammarians}.)

Source: \href{https://www.unicode.org/charts/PDF/U0980.pdf}{Unicode Bengali chart}.
\end{tcolorbox}
\section[chô]{\bn{চ} — the consonant cha}

\label{lesson:BN-W04-cha}



\begin{quote}
Write \textbf{\bn{দ}} and \textbf{\bn{ধ}}, and say which of the two lets the breath out.
\end{quote}

\subsection*{Script}
\begin{quote}
\bn{চ}
\end{quote}

\textbf{chô} — the \emph{ch} of English \emph{cheap}, with no puff of air after it.

The tongue is somewhere new. Not on the teeth, where \textbf{\bn{ত}} and \textbf{\bn{দ}} are made, and not at the back, where \textbf{\bn{ক}} is: the body of the tongue comes up flat against the \textbf{hard roof of the mouth}, a little behind the ridge. English spells that place with two letters, \emph{ch}, because its alphabet has no single sign for it. Bengali has one letter, and — following the square this chapter has been filling — a breathy partner and two voiced cells to go with it.

That is the fourth place of articulation this book has named: teeth, roof, throat, and lips. The list is not long. \textbf{Everything the script writes is one of those places with the voice and breath switches set.}

With the mātrā:

\begin{quote}
\textbf{\bn{চ}} + \textbf{\bn{া}} $\to$ \textbf{\bn{চা}}
\end{quote}

That is already a word, and it gets its own lesson next.

\subsection*{Your turn}
\begin{itemize}
\raggedright
  \item \emph{Write it:} \bn{চ}
  \item \emph{Write it:} \bn{চা}
  \item \emph{Read it:} \bn{কেমন} \bn{ভালো}
\end{itemize}

\begin{tcolorbox}[breakable,colback=teal!4,colframe=teal!35!black,title={Before you move on}]
Where in the mouth is \textbf{\bn{চ}} made? (Against the \textbf{roof}, behind the ridge.) How many letters does English need for that sound? (\textbf{Two}.) And Bengali? (\textbf{One}.)

Source: \href{https://www.unicode.org/charts/PDF/U0980.pdf}{Unicode Bengali chart}.
\end{tcolorbox}
\section[chhô]{\bn{ছ} — the breathy one on the roof}

\label{lesson:BN-W04-chha}



\begin{quote}
Write \textbf{\bn{চ}} and say where in the mouth it is made. Then write \textbf{\bn{গ}}.
\end{quote}

\subsection*{Script}
\begin{quote}
\bn{ছ}
\end{quote}

\textbf{chhô} — \textbf{\bn{চ}} with the breath let out behind it, the way \textbf{\bn{খ}} is \textbf{\bn{ক}} with the breath let out behind it. Same place on the roof of the mouth, same tongue, one switch flipped.

The pair is worth hearing rather than reading about. Hold a hand in front of your mouth. Say \emph{chô}: the hand feels almost nothing. Say \emph{chhô}: the hand is pushed. English has no contrast here at all — the \emph{ch} of \emph{church} sits somewhere between the two — so the ear has to be told what to listen for before it can hear it.

With the mātrā:

\begin{quote}
\textbf{\bn{ছ}} + \textbf{\bn{া}} $\to$ \textbf{\bn{ছা}}
\end{quote}

This letter has been under your nose since the first chapter. \textbf{\bn{আচ্ছা}} (\emph{āchchhā}, \textquotedblleft{}okay, I see\textquotedblright{}) is spelled with \textbf{\bn{চ}} and \textbf{\bn{ছ}} stacked, one on the other: the plain one, then the breathy one, said as a single long push. That is why the romanization has always carried that awkward \emph{chchh} — it was reporting two letters honestly.

\subsection*{Your turn}
\begin{itemize}
\raggedright
  \item \emph{Write it:} \bn{ছ}
  \item \emph{Write it:} \bn{চ} \bn{ছ}
  \item \emph{Write it:} \bn{ক} \bn{খ} / \bn{চ} \bn{ছ}
  \item \emph{Say it:} \emph{chô} then \emph{chhô}, hand in front of the mouth
\end{itemize}

\begin{tcolorbox}[breakable,colback=teal!4,colframe=teal!35!black,title={Before you move on}]
What separates \textbf{\bn{ছ}} from \textbf{\bn{চ}}? (\textbf{The breath}.) Which pair from the throat works the same way? (\textbf{\bn{ক}} and \textbf{\bn{খ}}.) And which word from the first chapter stacks these two letters? (\textbf{\bn{আচ্ছা}}.)

Source: \href{https://www.unicode.org/charts/PDF/U0980.pdf}{Unicode Bengali chart}.
\end{tcolorbox}
\section[āchhi]{\bn{আছি} — \textquotedblleft{}I am,\textquotedblright{} read}

\label{lesson:BN-W04-achhi-read}



\begin{quote}
Write \textbf{\bn{ছ}}. Then read \textbf{\bn{আমি}}.
\end{quote}

\subsection*{Script}
\begin{quote}
\textbf{\bn{আ}} + \textbf{\bn{ছি}} $\to$ \textbf{\bn{আছি}}
\end{quote}

\begin{itemize}
\raggedright
  \item \textbf{\bn{আ}} — the vowel standing on its own, because the word opens on it
  \item \textbf{\bn{ছ}} with \textbf{\bn{ি}} in front of it — \emph{chhi}
\end{itemize}

\textbf{āchhi} — \emph{\textquotedblleft{}{[}I{]} am {[}here{]}.\textquotedblright{}} It is the second half of the exchange you have been having since the responding chapter: \textbf{\bn{তুমি} \bn{কেমন} \bn{আছো}?} and the answer \textbf{\bn{আমি} \bn{ভালো} \bn{আছি}}.

Look at what the last sign is doing. \textbf{\bn{ি}} is not decoration and it is not part of the stem: it is the \textbf{I} of the sentence. Bengali packs the person into the verb's ending, so \emph{āchhi} means \textquotedblleft{}I am\textquotedblright{} with no pronoun needed at all, and \emph{āchho} means \textquotedblleft{}you are\textquotedblright{} by swapping that one sign for another. The \textbf{\bn{আমি}} in front is emphasis, not grammar.

That is a genuinely large idea arriving in a three-piece word, and it arrives by eye rather than by rule. You have been saying both endings for chapters. Seeing them written next to each other is what makes them a \emph{pattern} rather than two memorised noises.

Now the reply, whole:

\begin{quote}
\textbf{\bn{আমি} \bn{ভালো} \bn{আছি}}
\end{quote}

Ten pieces, three words, and not one new shape in any of them.

\subsection*{Your turn}
\begin{itemize}
\raggedright
  \item \emph{Read it:} \bn{আছি}
  \item \emph{Read it:} \bn{আমি} \bn{ভালো} \bn{আছি}
  \item \emph{Write it:} \bn{আছি}
  \item \emph{Say it:} the whole reply, out loud
\end{itemize}

\begin{tcolorbox}[breakable,colback=teal!4,colframe=teal!35!black,title={Before you move on}]
Which piece of \textbf{\bn{আছি}} carries the \textquotedblleft{}I\textquotedblright{}? (\textbf{The \bn{ি} at the end}.) So does the sentence need \textbf{\bn{আমি}}? (\textbf{No} — it is emphasis.) And how many new shapes were in \textbf{\bn{আমি} \bn{ভালো} \bn{আছি}}? (\textbf{None}.)

Source: \href{https://www.unicode.org/charts/PDF/U0980.pdf}{Unicode Bengali chart}.
\end{tcolorbox}
\section[jô]{\bn{জ} — the consonant ja, and the grid closed}

\label{lesson:BN-W04-ja}



\begin{quote}
Write \textbf{\bn{চ}} and \textbf{\bn{ছ}}, and say what separates them.
\end{quote}

\subsection*{Script}
\begin{quote}
\bn{জ}
\end{quote}

\textbf{jô} — the \emph{j} of English \emph{jam}. Same place as \textbf{\bn{চ}}, the flat of the tongue on the roof of the mouth, with the voice on.

The third cell of the roof, and by now the move should be predictable: meet a consonant, ask where it is made and whether the voice and the breath are running, and the letter has a family rather than being a shape on its own.

This letter carries a root worth knowing. Sanskrit \emph{jñā}, \textquotedblleft{}to know,\textquotedblright{} is spelled with it, and the verbs chapter followed that root out to English \textbf{know} and to the \emph{k} that English still writes and no longer says. The Bengali verb for knowing facts opens on this shape.

With the mātrā:

\begin{quote}
\textbf{\bn{জ}} + \textbf{\bn{া}} $\to$ \textbf{\bn{জা}}
\end{quote}

That is the last consonant of the two script chapters that built this grid, and it closes a table rather than lengthening a list. Every letter in it was met as a cell — a place in the mouth, with the voice and the breath switched on or off. Say them back by place rather than by shape:

\begin{itemize}
\raggedright
  \item \textbf{at the lips}: \textbf{\bn{প}} voice off · \textbf{\bn{ব}} voice on · \textbf{\bn{ভ}} voice on with breath
  \item \textbf{on the teeth}: \textbf{\bn{ত}} voice off · \textbf{\bn{দ}} voice on · \textbf{\bn{ধ}} voice on with breath
  \item \textbf{on the roof}: \textbf{\bn{চ}} voice off · \textbf{\bn{ছ}} voice off with breath · \textbf{\bn{জ}} voice on
  \item \textbf{at the throat}: \textbf{\bn{ক}} voice off · \textbf{\bn{খ}} voice off with breath · \textbf{\bn{গ}} voice on
\end{itemize}

Twelve cells, four rows, and not one of them learned as a picture. The cells nobody has filled are letters this book has not needed yet, not gaps in the script. A reader who can recite that table by place is filling in a grid; a reader holding twelve shapes is memorising twelve shapes. That difference is the whole reason these chapters kept asking \emph{where} and \emph{with what breath} instead of \emph{what does this look like}.

\subsection*{Your turn}
\begin{itemize}
\raggedright
  \item \emph{Write it:} \bn{জ}
  \item \emph{Write it:} \bn{চ} \bn{ছ} \bn{জ}
  \item \emph{Write it:} \bn{প} \bn{ব} \bn{ভ} / \bn{ত} \bn{দ} \bn{ধ} / \bn{চ} \bn{ছ} \bn{জ} / \bn{ক} \bn{খ} \bn{গ}
  \item \emph{Read it:} \bn{আছি}
  \item \emph{Read it:} \bn{আবার} \bn{ভালো}
\end{itemize}

\begin{tcolorbox}[breakable,colback=teal!4,colframe=teal!35!black,title={Before you move on}]
What separates \textbf{\bn{জ}} from \textbf{\bn{চ}}? (\textbf{The voice}.) How many cells does the grid now hold? (\textbf{Twelve} — four rows of three.) And what question does the grid ask of a new letter — what it looks like, or where it is made? (\textbf{Where it is made}, and with what breath.)

Source: \href{https://www.unicode.org/charts/PDF/U0980.pdf}{Unicode Bengali chart}.
\end{tcolorbox}
